\originalsection{作业 3: 圆周运动}
\begin{exercise}
原题 1. 光滑竖直圆环半径 $R$, 绕竖直直径以角速度 $\omega$ 转动. 环上小珠在相对最低竖直方向角 $\theta$ 处相对环静止. (a) 求所需 $\omega$. (b) 低于何阈值 $\omega_0$ 时仅底端平衡?
\end{exercise}
\begin{solution}
珠沿环切向的有效力为 $-mg\sin\theta+m\omega^2R\sin\theta\cos\theta$. 非底端平衡给 $\omega^2=g/(R\cos\theta)$, 故 $\omega=\sqrt{g/(R\cos\theta)}$. 存在要求 $\cos\theta\le1$, 所以 $\omega_0=\sqrt{g/R}$. 严格说环顶也形式上是平衡但不稳定, 原题“only balanced position”应理解为稳定平衡; 不能在完整平衡集合中删去顶端.
\end{solution}
\begin{exercise}
原题 2. 汽车质量 $m$, 转弯半径 $R$, 路面横向倾角 $\beta$, 静摩擦系数 $\mu_s$. 求无滑速度范围 $v_{\min}\le v\le v_{\max}$.
\end{exercise}
\begin{solution}
高速时摩擦沿坡向下, $N\sin\beta+\mu_sN\cos\beta=mv^2/R$, $N\cos\beta-\mu_sN\sin\beta=mg$. 故 $v_{\max}^2=Rg(\sin\beta+\mu_s\cos\beta)/(\cos\beta-\mu_s\sin\beta)$. 低速极限摩擦向上, $v_{\min}^2=Rg(\sin\beta-\mu_s\cos\beta)/(\cos\beta+\mu_s\sin\beta)$. 分子非正时最低速度为零. 高速式分母非正时理想模型不给有限上界, 须另查接触和路面约束. 通常角度范围为 $0\le\beta<\pi/2$.
\end{solution}
\begin{exercise}
原题 3. 长 $l$ 的绳上小球作圆锥摆, 绳与竖直夹角 $\beta$, 轨道离地高度 $h$. 忽略阻力. (a) 求角速度. (b) 球经过 $x$ 轴时断绳, 求落地前水平位移大小 $d$.
\end{exercise}
\begin{solution}
(a) 第\ref{ch:newton}章得 $\omega=\sqrt{g/(l\cos\beta)}$. (b) 初始水平速率 $v=\omega l\sin\beta$, 飞行时间 $\sqrt{2h/g}$, 故 $d=\sin\beta\sqrt{2hl/\cos\beta}$. 断绳后的方向沿原圆的切线, 不是半径方向.
\end{solution}
\begin{exercise}
原题 4. 两盒绕轴以角速度 $\omega$ 转动, 盒 1 质量 $m_1$ 距轴 $d$, 盒 2 质量 $m_2$ 距轴 $2d$. 绳 A 连接轴和盒 1, 绳 B 连接盒 1 和盒 2. 忽略绳质量与重力. (a) 求 $T_B$. (b) 求 $T_A$.
\end{exercise}
\begin{solution}
盒 2 只有 B 的向心力, $T_B=2m_2d\omega^2$. 盒 1 受 A 向内、B 向外, $T_A-T_B=m_1d\omega^2$, 故 $T_A=(m_1+2m_2)d\omega^2$.
\end{solution}
\begin{exercise}
原题 5. 两个同质量卫星处于不同高度的地球圆轨道. (a) 哪个速度大? (b) 哪个周期长?
\end{exercise}
\begin{solution}
向心关系给 $v=\sqrt{GM_E/r}$, $T=2\pi\sqrt{r^3/(GM_E)}$. 较低卫星更快, 较高卫星周期更长. 半径从地心量而非从地面量.
\end{solution}
\begin{exercise}
原题 6. 质量 $m$ 的硬币在圆盘距轴 $d$ 处, 静摩擦系数 $\mu_s$. 圆盘从静止以恒定角加速度 $\alpha$ 启动. (a) 无滑前静摩擦大小随时间如何变化? (b) 开始滑动时角速度为多少?
\end{exercise}
\begin{solution}
径向加速度 $-d\alpha^2t^2$, 切向 $d\alpha$. 所以 $f_s=md\sqrt{\alpha^2+\alpha^4t^4}$. 达到 $\mu_smg$ 时 $\omega_*=[(\mu_sg/d)^2-\alpha^2]^{1/4}$. 若 $\alpha>\mu_sg/d$, 立即滑动; 等号时极限启动瞬间即达阈值.
\end{solution}
